\documentclass{article}

\begin{document}
    \begin{center}
        \Large \textbf{Challenge 2: Persistenz} \\ [6pt]
        \normalsize
        Autor: Maximilian Jakob Maag B.Sc. \\
        Version: 1.0
    \end{center}
    \section{Intro}
    Wenn Terraform ein Element – z.B. eine Firewall, VPS o.Ä. – provisioniert, das bereits existiert, wird das bestehende Element zerstört und ein neues
    Element provisioniert. Das kann dazu führen, dass Produktionsdaten oder Konfigurationen verloren gehen, wenn man nicht aufpasst.
    \section{Aufgaben}
In dieser Übung sollen Sie einige VPS‑Instanzen deployen und Konfigurationen vornehmen. Danach werden Sie diese ändern.
    Ihre Aufgabe ist es, dafür zu sorgen, dass keine Daten verloren gehen. \\ \\
    Um Anwendungsdaten zu simulieren, können Sie beispielsweise einen Mailserver wie Mailcow deployen.
    \subsection{Aufgabe 1}
    Suchen Sie sich eine Anwendung aus, welche Sie deployen wollen. Als Beispiel können Sie u.a. Mailcow verwenden. \\
    Deployen Sie zunächst ein VPS, das so konfiguriert wurde das sie in der Lage sind sich mittels SSH zu verbinden. \\ 
    Verbinden Sie sich mit dem VPS und deployen Sie auf dem VPS mailcow mittels Docker Compose. \\ \\
    Ändern Sie eine Kleinigkeit an Ihrer Terraform‑Konfiguration, z.B. die RAM-Größe Ihres VPS von 4 auf 8 GB.
    Deployen Sie Ihre Änderung mittels \texttt{terraform apply}. \\ \\
    Beobachten Sie das Verhalten des Mailservers. Was hat sich verändert? Ist der Server noch erreichbar? Welche IP-Adresse hat Ihr VPS jetzt?
    \subsection{Aufgabe 2}
    Ihr Mailserver sollte nun nicht mehr existieren. Erläutern Sie, warum das der Fall sein könnte. \\
    Ändern Sie Ihre Konfiguration. Deployen Sie neben Ihrem VPS ein Volume mit einer Größe von 10 GB und verbinden Sie dieses mit Ihrem VPS.
    Mounten Sie das Volume unter /mnt/mail-data.
    Überlegen Sie, wie sie auf diesem Volume die Konfiguration von Mailcow und ein Backup von Mailcow sinnvoll speichern  können.
    Automatisieren Sie anschließend das Deployment eines VPS den Mount des Volumes, das Starten von Mailcow und das Einspielen des Backups auf der neuen Mailcow Instanz.
    \subsection{Aufgabe 3}
    Beobachten Sie das Verhalten des E-Mailservers erneut. Was stellen Sie fest, wenn Sie die Konfiguration Ihres VPS ändern? Was stellen Sie fest wenn Sie die Konfiguration des Volumes ändern?
    Welches Desaster könnte entstehen und wie könnten Sie dieses beheben?
    \subsection{Aufgabe 4}
    Führen Sie das in Aufgabe 3 von Ihnen beschriebene Desaster herbei und beheben Sie es.  \\ \\
    Zum Schluss führen Sie nach dem Abschluss der Übung terraform destroy aus.
\end{document}